\begin{table}[H]
\caption{Configuraciones del segundo nivel de la validación empírica}
\label{tab:configuraciones}
\setstretch{1}
\renewcommand{\arraystretch}{1.3}
\small
\begin{tabularx}{\textwidth}{|L{1.3cm}|L{4.6cm}|Y|}
\hline
\rowcolor{gray!25}
\textbf{Config.} & \textbf{Composición} & \textbf{Propósito} \\
\hline
$C_1$ & SLM sin recuperación & Línea base generativa: desempeño del modelo con su solo conocimiento paramétrico \\
\hline
$C_2$ & SLM + RAG convencional & Aporte de la recuperación sobre el corpus empresarial \\
\hline
$C_3$ & SLM + RAG + crítico paramétrico (Self-RAG) & Aporte de la autorreflexión cuya especificidad de dominio reside en los pesos (\cite{Asai2024_Selfrag}) \\
\hline
$C_4$ & SLM + RAG + Critic-EmpRAG1.0 & Sistema completo: crítico híbrido con juicio anclado en la evidencia recuperada \\
\hline
$C_{4a}$ & $C_4$ sin ciclo correctivo & Ablación del refinamiento iterativo \\
\hline
$C_{4b}$ & $C_4$ sin evidencia externa para el crítico & Ablación del componente no paramétrico del crítico \\
\hline
\end{tabularx}

\vspace{0.2cm}
\noindent{\small \textit{Nota.} Todas las configuraciones comparten el modelo base y el corpus documental; de $C_2$ en adelante comparten además el mecanismo de recuperación y los parámetros de segmentación e indexación. SLM: modelo de lenguaje pequeño; RAG: generación aumentada por recuperación.}
\end{table}
